\SourceLeaf{270}%
und folglich, wenn man mithilfe von 5 entweder \(\mathrm{d}U\) oder \(\mathrm{d}w\) eliminiert,
\[
\frac{\mathrm{d}\log m}{\mathrm{d}w}=\frac{(1-ee)(\alpha\cos U-\cos w)}{(1-ee\cos w^2)\sin w} \quad .\ .\ .\ .\ .\ .\ .\ .\ .\ .\ .\ .\quad(6)
\]
\[
\frac{\mathrm{d}\log m}{\mathrm{d}U}=\operatorname{cotg}U-\frac{\cos w}{\alpha\sin U}=\frac{\alpha\cos U-\cos w}{\alpha\sin U} \quad .\ .\ .\ .\ .\ .\ .\quad(7)
\]

Differenziert man Gleichung 7 ein weiteres Mal, so erhält man
\[
\begin{aligned}
\frac{\mathrm{dd}\log m}{\mathrm{d}U^2}&=-\frac{1}{\sin U^2}+\frac{\cos U\cos w}{\alpha\sin U^2}+\frac{\sin w}{\alpha\sin U}\mathord{.}\frac{\mathrm{d}w}{\mathrm{d}U}\\
&=-\frac{1}{\sin U^2}+\frac{\cos U\cos w}{\alpha\sin U^2}+\frac{(1-ee\cos w^2)\sin w^2}{\alpha\alpha(1-ee)\sin U^2} \quad .\ .\ .\ .\ .\ .\quad(8)
\end{aligned}
\]

Soll \(m\) für eine bestimmte Breite (Normalbreite) gleich 1 sein und für andere Breiten nur um Größen dritter Ordnung von 1 abweichen, wobei die Breitenunterschiede als Größen erster Ordnung betrachtet werden, so bezeichnen wir die Normalbreite auf dem Ellipsoid mit \(P\) und die entsprechende Breite auf der Kugel mit \(Q\). Dann müssen für \(w=90^0-P\), \(U=90^0-Q\) nach den Gleichungen 4, 7 und 8 die folgenden Beziehungen gelten.
\[
A=\frac{a\cos P}{\alpha\cos Q\BareRoot{(1-ee\sin P^2)}} \quad .\ .\ .\ .\ .\ .\ .\ .\ .\ .\ .\ .\ .\ .\quad(9)
\]
\[
\alpha\sin Q=\sin P \quad .\ .\ .\ .\ .\ .\ .\ .\ .\ .\ .\ .\ .\ .\ .\ .\ .\ .\ .\quad(10)
\]
\[
0=1-\frac{\sin P\sin Q}{\alpha}-\frac{(1-ee\sin P^2)\cos P^2}{\alpha\alpha(1-ee)}
\]
Setzt man in der letzten Gleichung für \(\sin Q\) seinen Wert aus 10 ein, so ergibt sich
\[
\alpha\alpha=1+\frac{ee\cos P^4}{1-ee} \quad .\ .\ .\ .\ .\ .\ .\ .\ .\ .\ .\ .\ .\ .\ .\ .\quad(11)
\]

Diese Gleichung liefert also \(\alpha\), sobald für \(P\) ein bestimmter Wert gewählt ist. Anschließend lassen sich \(Q\) aus Gleichung 10 und \(A\) aus Gleichung 9 bestimmen. Schließlich ergibt sich \(k\) durch Einsetzen von \(w=90^0-P\), \(U=90^0-Q\) in die allgemeine Gleichung 3, nämlich
\[
k=\frac{\operatorname{tang}(45^0+\frac{1}{2}P)^\alpha}{\operatorname{tang}(45^0+\frac{1}{2}Q)}\mathord{.}(\tfrac{1-e\sin P}{1+e\sin P})^{\frac{1}{2}\alpha e} \quad .\ .\ .\ .\ .\ .\ .\ .\quad(12)
\]

\Article{4}
\indent Die Konstanten \(A,\alpha,k\) und die Normalbreite auf der Kugel \(Q\) wird man aus \(P\) und \(e\) mit besonderer Sorgfalt und Genauigkeit berechnen wollen, da all diese Größen als Grundlage für
